\begin{proof}
    \textbf{($\subseteq$)} Let $p \in E \setminus E'$ be given.
    Since $p \notin E'$, there exists some $r \in \R^{+}$ such that
    \[
    E \cap \widetilde{N}_{r}(p) = \emptyset \implies \paren*{E \cap N_{r}(p)} \setminus \set*{p} = \emptyset.
    \]
    Since $p \in E$ and $p \in N_{r}(p)$, we have
    \[
    E \cap N_{r}(p) = \set*{p},
    \]
    which implies that $p$ is an isolated point of $E$.
    \\[6mm]
    \textbf{($\supseteq$)} Let $p$ be an isolated point of $E$. Then there exists some $r \in \R^{+}$ such that
    \[
    E \cap N_{r}(p) = \set*{p}.
    \]
    Then we have
    \begin{itemize}
        \item $E \cap \widetilde{N}_{r}(p) = \emptyset \implies p \notin E'$.
        \item $p \in E$.
    \end{itemize}
    Therefore, we have $p \in E \setminus E'$.
\end{proof}

\bigskip

\begin{definition} \label{def:Boundary}
    Let $(X, d)$ be a metric space and $E \subseteq X$. The \textbf{boundary} of $E$,
    which is denoted by $\partial E$, is defined by
    \[
    \partial E \coloneqq \overline{E} \cap \overline{E^{c}}.
    \]
\end{definition}

\bigskip

\begin{lemma} \label{lem:DeletedNeighborhoodInRnIsNonEmpty}
    Let $(\R^{n}, \norm*{\cdot})$ be given for some $n \in \N$. Then we have
    \[
    \widetilde{N}_r(\bm{p}) \neq \emptyset \quad \text{for all } \bm{p} \in \R^{n} \text{ and } r \in \R^{+}.
    \]
\end{lemma}

\begin{proof}
    Let $\bm{p} = (p_{1}, p_{2}, \dots, p_{n}) \in \R^{n}$ and $r \in \R^{+}$ be given. Then we have
    \[
    d\paren*{\bm{p} + \frac{r}{2} \bm{e_{1}}, \bm{p}} = \norm*{\paren*{\bm{p} + \frac{r}{2} \bm{e_{1}}}
    - \bm{p}} = \norm*{\frac{r}{2} \bm{e_{1}}} = \frac{r}{2} \norm*{\bm{e_{1}}} = \frac{r}{2} < r.
    \]
    where $\bm{e_{1}} = (1, 0, 0, \dots, 0) \in \R^{n}$. Therefore, we have
    \[
    \bm{p} \neq \paren*{\bm{p} + \frac{r}{2} \bm{e_{1}}} \in \widetilde{N}_r(\bm{p}) \implies \widetilde{N}_r(\bm{p}) \neq \emptyset.
    \]
\end{proof}

\bigskip

\begin{theorem} \label{thm:PropertiesOfBoundaryInRn}
    Let $(\R^{n}, \norm*{\cdot})$ be given for some $n \in \N$ and $E \subseteq \R^{n}$. Then we have
    \[
    \partial E = \paren*{E \setminus E'} \cup \paren*{E^{c} \setminus (E^{c})'} \cup \paren*{E' \cap (E^{c})'},
    \]
    which states that the boundary of $E$ consists the following three disjoint sets:
    \begin{itemize}
        \item The set of all isolated points of $E$.
        \item The set of all isolated points of $E^{c}$.
        \item The set of all limit points of both $E$ and $E^{c}$.
    \end{itemize}
\end{theorem}

\begin{proof}
    By the definition of the boundary, we have
    \begin{align*}
        \partial E &= \overline{E} \cap \overline{E^{c}} \\
        &= \paren*{E \cup E'} \cap \paren*{E^{c} \cup (E^{c})'} \\
        &= \paren*{E \cap E^{c}} \cup \paren*{E \cap (E^{c})'}
        \cup \paren*{E' \cap E^{c}} \cup \paren*{E' \cap (E^{c})'}.
    \end{align*}
    Consider the following equivalences:
    \begin{itemize}
        \item $E = \paren*{E \setminus E'} \cup \paren*{E \cap E'}$.
        \item $E^{c} = \paren*{E^{c} \setminus (E^{c})'} \cup \paren*{E^{c} \cap (E^{c})'}$.
    \end{itemize}
    Then we have the following equalities:
    \begin{itemize}
        \item $E \cap (E^{c})' = \paren*{\paren*{E \setminus E'} \cap (E^{c})'}
        \cup \paren*{\paren*{E \cap E'} \cap (E^{c})'}$.
        \item $E' \cap E^{c} = \paren*{E' \cap \paren*{E^{c} \setminus (E^{c})'}}
        \cup \paren*{E' \cap \paren*{E^{c} \cap (E^{c})'}}$.
    \end{itemize}
    Consider the last two terms of the above equalities:
    \begin{itemize}
        \item $\paren*{\paren*{E \cap E'} \cap (E^{c})'} \subseteq E' \cap (E^{c})'$.
        \item $\paren*{E' \cap \paren*{E^{c} \cap (E^{c})'}} \subseteq E' \cap (E^{c})'$.
    \end{itemize}
    Then we have
    \begin{align*}
        \partial E &= \emptyset \cup \paren*{\paren*{E \setminus E'} \cap (E^{c})'}
        \cup \paren*{E' \cap \paren*{E^{c} \setminus (E^{c})'}} \cup \paren*{E' \cap (E^{c})'} \\
        &= \paren*{\paren*{E \setminus E'} \cap (E^{c})'}
        \cup \paren*{E' \cap \paren*{E^{c} \setminus (E^{c})'}} \cup \paren*{E' \cap (E^{c})'}
    \end{align*}
    For any $\bm{p} \in E \setminus E'$, there exists some $r \in \R^{+}$ such that
    \[
    E \cap \widetilde{N}_{r}(\bm{p}) = \emptyset \implies \widetilde{N}_{r}(\bm{p}) \subseteq E^{c}.
    \]
    By \Cref{lem:DeletedNeighborhoodInRnIsNonEmpty}, we have
    \[
    \widetilde{N}_{s}(\bm{p}) \cap \widetilde{N}_{r}(\bm{p}) = \widetilde{N}_{\min(s,r)}(\bm{p}) \neq \emptyset\quad \text{for all } s \in \R^{+}.
    \]
    Then, for any $s \in \R^{+}$, we have
    \begin{align*}
        \widetilde{N}_{s}(\bm{p}) \cap \widetilde{N}_{r}(\bm{p}) \neq \emptyset &\implies \widetilde{N}_{s}(\bm{p}) \cap
        \paren*{\widetilde{N}_{r}(\bm{p}) \cap E^{c}} \neq \emptyset \\
        &\implies \emptyset \neq \paren*{\widetilde{N}_{s}(\bm{p}) \cap E^{c}} \cap \widetilde{N}_{r}(\bm{p})
        \subseteq \widetilde{N}_{s}(\bm{p}) \cap E^{c} \\
        &\implies E^{c} \cap \widetilde{N}_{s}(\bm{p}) \neq \emptyset,
    \end{align*}
    which implies that $\bm{p} \in (E^{c})'$. Then we have
    \[
    E \setminus E' \subseteq (E^{c})' \implies \paren*{E \setminus E'} \cap (E^{c})' = E \setminus E'.
    \]
    Similarly, we can show that
    \[
    E^{c} \setminus (E^{c})' \subseteq E' \implies E' \cap \paren*{E^{c} \setminus (E^{c})'} = E^{c} \setminus (E^{c})'.
    \]
    Therefore, we have
    \[
    \partial E = \paren*{E \setminus E'} \cup \paren*{E^{c} \setminus (E^{c})'} \cup \paren*{E' \cap (E^{c})'}.
    \]
\end{proof}

\bigskip

\begin{remark}
    Let $(X, d)$ be a metric space and $Y \subseteq X$. Then, $(Y, d|_{Y \times Y})$ is also a metric space, since
    \[
    d|_{Y \times Y}: Y \times Y \to \R
    \]
    satisfies the properties of a metric. We call $(Y, d|_{Y \times Y})$ a \textbf{subspace} of $(X, d)$. For the
    sake of simplicity, we will denote $d|_{Y \times Y}$ by $d$ and write $(Y, d)$ instead of $(Y, d|_{Y \times Y})$.
\end{remark}

\bigskip

\begin{definition} \label{def:OpenSetRelativeToSubspace}
    Let $(X, d)$ be a metric space and $E \subseteq Y \subseteq X$. $E$ is called \textbf{open relative} to $Y$
    if $E$ is open in the subspace $(Y, d)$. That is, for any $p \in E$, there exists some $r \in \R^{+}$ such that 
    \[
    Y \cap N_{r}(p) \subseteq E.
    \]
\end{definition}

\bigskip

\begin{theorem} \label{thm:OpenSetRelativeToSubspace}
    Let $(X, d)$ be a metric space and $E \subseteq Y \subseteq X$. Then $E$ is open relative to $Y$
    if and only if there exists an open set $V$ in $X$ such that $E = V \cap Y$.
\end{theorem}

\begin{proof}
    \textbf{($\implies$)} Let $E$ be an open set relative to $Y$.
    Then, for each $p \in E$, there exists some $r_{p} \in \R^{+}$ such that
    \[
    Y \cap N_{r_{p}}(p) \subseteq E.
    \]
    Let $V$ be a set defined by
    \[
    V \coloneqq \bigcup_{p \in E} N_{r_{p}}(p).
    \]
    By \Cref{thm:NeighborhoodIsOpen,thm:ArbitraryUnionOfOpenSetsIsOpen}, $V$ is open in $X$.
    Since $p \in N_{r_{p}}(p) \subseteq V$ for each $p \in E$, we have
    \begin{itemize}
        \item $E \subseteq Y$ and $E \subseteq V \implies E \subseteq V \cap Y$.
        \item $\displaystyle \bigcup_{p \in E} \paren*{Y \cap N_{r_{p}}(p)} \subseteq E \implies Y \cap
        \bigcup_{p \in E} N_{r_{p}}(p) \subseteq E \implies Y \cap V \subseteq E$.
    \end{itemize}
    Therefore, we have $E = V \cap Y$.
    \\[6mm]
    \textbf{($\impliedby$)} Let $V$ be an open set in $X$ such that $E = V \cap Y$.
    Then, for each $p \in E \subseteq V$, there exists some $r_{p} \in \R^{+}$ such that
    \[
    N_{r_{p}}(p) \subseteq V \implies Y \cap N_{r_{p}}(p) \subseteq V \cap Y = E.
    \]
\end{proof}

\end{section}